

\section{Twisted endomorphism bicategories}\label{sec:remnant}
In this section we verify the existence of the endomorphism bicategory $\Lambda \mathcal{B}$ as in the following statement.

\restateLambdabicat*

We note here that the assumption that some of the 1-cells are dualizable and some 2-cells are isomorphisms are forward looking assumptions for \cref{sec:traces}.  They are not required to define a bicategory and lead to more work in this section.


We start by verifying that $\Lambda \mathcal{B}(P,Q)$ is a category for endomorphism 1-cells $P$ and $Q$ of $\mathcal{B}$.

\begin{lem}\label{LambdaBPQ_Cat}
    For endomorphism 1-cells $P$ and $Q$  in $\B$, $\Lambda \B(P, Q)$ is a category.
\end{lem}

\begin{proof}
The composition of morphisms in $\Lambda \B(P,Q)$, denoted $\Box$, is given by composing the associated 2-cells from $\B$.  The required diagrams commute since they correspond to vertically stacked commutative diagrams. 
\[\xymatrix{P\odot L \ar@{=>}[r]^\delta \ar@{=>}[d]_{\id\odot\phi} & L\odot Q \ar@{=>}[d]^{\phi\odot\id}\\
P\odot M \ar@{=>}[d]_{\id\odot\psi} \ar@{=>}[r]^\beta & M\odot Q \ar@{=>}[d]^{\psi\odot\id}\\
P\odot N \ar@{=>}[r]^\gamma & N\odot Q
}.\] 


       \underline{Associativity}: Composition of 2-cells in $\Lambda \B$ is composition in $\B$. 
       Since composition of 2-cells in $\B$ is associative,  composition of morphisms in $\Lambda \B(P, Q)$ is associative. 
        
        \underline{Identity Morphism}: Let $(M, \beta): P \to Q$ be an object in $\Lambda \B(P, Q)$. Then $\id_M$ defines an endomorphism of $(M,\beta)$ since the following diagram commutes.  
        \[\xymatrix{P\odot M \ar@{=>}[d]_{\id_P\odot\id_M} \ar@{=>}[r]^\beta & M\odot Q \ar@{=>}[d]^{\id_M\odot \id_Q}\\
        P\odot M \ar@{=>}[r]_\beta & M\odot Q
        }\] 
        To see that 
        \[\ltwo{\id}_M\Box \ltwo{\psi} = \ltwo{\psi}\quad \text{and}\quad  \ltwo{\psi}\Box \ltwo{\id}_N = \ltwo{\psi}\] 
        for a 2-cell $\ltwo{\psi}: (M, \beta) \Rightarrow (N, \gamma)$ in $\Lambda \B(P, Q)$, 
        it is enough to note that the composition $\Box$ is composition of 2-cells in $\B$.

Since the structure maps are inherited from $\mathcal{B}$ the required diagrams commute. 
\end{proof}

Given objects $(M, \beta)\in \Lambda\B( P,Q)$ and $(N, \gamma)\in \Lambda \B( Q,T)$ in $\Lambda \B$, let 
$\boxmaps{\beta}{\gamma}$ 
be the composite 
% https://q.uiver.app/#q=WzAsMyxbMCwwLCJQXFxvZG90IE1cXG9kb3QgTiJdLFsxLDEsIk0gXFxvZG90IFFcXG9kb3QgTiJdLFsyLDAsIk0gXFxvZG90IE5cXG9kb3QgVCJdLFswLDIsIlxcYm94bWFwc3tcXGJldGF9e1xcZ2FtbWF9IiwwLHsibGV2ZWwiOjJ9XSxbMCwxLCJ7XFxiZXRhXFxvZG90IFxcbWF0aHJte2lkfV9OfSIsMix7ImxldmVsIjoyfV0sWzEsMiwie1xcbWF0aHJte2lkfV9NXFxvZG90IFxcZ2FtbWF9IiwyLHsibGV2ZWwiOjJ9XV0=&macro_url=https%3A%2F%2Fwww.overleaf.com%2Fproject%2F64c9c0fcea71c453e68a663b
\[\begin{tikzcd}
	{P\odot M\odot N} && {M \odot N\odot T} \\
	& {M \odot Q\odot N}
	\arrow["{\boxmaps{\beta}{\gamma}}", Rightarrow, from=1-1, to=1-3]
	\arrow["{{\beta\odot \mathrm{id}_N}}"', Rightarrow, from=1-1, to=2-2]
	\arrow["{{\mathrm{id}_M\odot \gamma}}"', Rightarrow, from=2-2, to=1-3]
\end{tikzcd}\]
Then the  bicategorical composition in $\Lambda \B$ of $(M,\beta)$ and $(N,\gamma)$, denoted  
\[(M, \beta) \boxdot (N, \gamma),\]
is $(M\odot N, \boxmaps{\beta}{\gamma})$. 

\begin{notn}\label{notation_box_dot}
   We will abuse notation and  also use $(M, \beta) \boxdot (N, \gamma)$  to indicate  the  map $\boxmaps{\beta}{\gamma}$.
\end{notn}

\begin{lem}\label{LambdaB_assoc}
 For 1-cells $(L, \delta): P\to Q, (M, \beta): Q\to R$, and $(N, \gamma): R \to S$ in $\Lambda\B$, the associator $a_{L,M,N}\colon (L\odot M)\odot N\to L\odot (M\odot N)$ in $\B$ is a 2-cell 
\[(L, \delta)\boxdot\left[(M, \beta)\boxdot(N, \gamma)\right]
\to (\left[L, \delta)\boxdot(M, \beta)\right]\boxdot(N, \gamma)\] in $\Lambda\B$.
\end{lem}

   In particular, $a_{-,-,-}$ is an associator for $\Lambda\B$ satisfying the pentagon axiom.

\begin{proof}
    We need to show 
the diagram in \eqref{eq:assocator_1-cells} commutes. 
    \begin{equation}\label{eq:assocator_1-cells}
    \xymatrix@C=4cm{P\odot \left[L\odot (M\odot N)\right] \ar@{=>}[r]^{(L, \delta)\boxdot\left[(M, \beta)\boxdot(N, \gamma)\right]} \ar@{<=}[d]_{\id_P\odot a_{L, M, N}} &  \left[L\odot (M\odot N)\right]\odot S \ar@{<=}[d]^{a_{L, M, N}\odot\id_S}\\
    P\odot \left[(L\odot M)\odot N\right] \ar@{=>}[r]_{\left[(L, \delta)\boxdot(M, \beta)\right]\boxdot(N, \gamma)} & \left[(L\odot M)\odot N\right]\odot S}.
    \end{equation}
The top and bottom maps in \eqref{eq:assocator_1-cells} are  the 2-cells 
in \cref{eq:assoctivity_Lambda_B}.  Here we have suppressed the associators  \cite{bicat_coherence,MP:coherence}.  This has the additional benefit that it is easier to see that the diagram in  \eqref{eq:assocator_1-cells} commutes. 
\begin{figure}
     \centering
     \begin{subfigure}[b]{.45\textwidth}\label{lamdba_b_2_1}
% https://q.uiver.app/#q=WzAsNCxbMSwwLCJMXFxvZG90IE1cXG9kb3QgTlxcb2RvdCBTIl0sWzAsMCwiUFxcb2RvdCBMXFxvZG90IE1cXG9kb3QgTiAiXSxbMSwxLCJMXFxvZG90IE1cXG9kb3QgUlxcb2RvdCBOKSAiXSxbMCwxLCJMXFxvZG90IChRXFxvZG90IE0pXFxvZG90IE4iXSxbMSwyLCJ7KEwsXFxkZWx0YSlcXGJveGRvdChNLCBcXGJldGEpXFxvZG90XFxtYXRocm17aWR9X059IiwxLHsibGV2ZWwiOjJ9XSxbMiwwLCJ7XFxtYXRocm17aWR9X3tMXFxvZG90IE19XFxvZG90IFxcZ2FtbWF9IiwyLHsibGV2ZWwiOjJ9XSxbMSwzLCJ7XFxkZWx0YSBcXG9kb3QgXFxtYXRocm17aWR9X01cXG9kb3QgXFxtYXRocm17aWR9X059IiwyLHsibGV2ZWwiOjJ9XSxbMywyLCJ7XFxtYXRocm17aWR9X0xcXG9kb3QgXFxiZXRhXFxvZG90IFxcbWF0aHJte2lkfV9OfSIsMix7ImxldmVsIjoyfV0sWzEsMCwieyhMLCBcXGRlbHRhKVxcYm94ZG90KE0sIFxcYmV0YSlcXGJveGRvdChOLCBcXGdhbW1hKX0iLDAseyJsZXZlbCI6Mn1dXQ==
\adjustbox{scale=.9,center}{\begin{tikzcd}[column sep=1in,row sep=.5in]
	{P\odot L\odot M\odot N } & {L\odot M\odot N\odot S} \\
	{L\odot (Q\odot M)\odot N} & {L\odot M\odot R\odot N) }
	\arrow["{{(L, \delta)\boxdot(M, \beta)\boxdot(N, \gamma)}}", Rightarrow, from=1-1, to=1-2]
	\arrow["{{\delta \odot \mathrm{id}_M\odot \mathrm{id}_N}}"', Rightarrow, from=1-1, to=2-1]
	\arrow["{{(L,\delta)\boxdot(M, \beta)\odot\mathrm{id}_N}}"{description}, Rightarrow, from=1-1, to=2-2]
	\arrow["{{\mathrm{id}_L\odot \beta\odot \mathrm{id}_N}}"', Rightarrow, from=2-1, to=2-2]
	\arrow["{{\mathrm{id}_{L\odot M}\odot \gamma}}"', Rightarrow, from=2-2, to=1-2]
\end{tikzcd}}
\caption{$\left[(L, \delta)\boxdot(M, \beta)\right]\boxdot(N, \gamma)$}
\end{subfigure}
\hfill 
     \begin{subfigure}[b]{0.45\textwidth}\label{lamdba_b_1_2}
% https://q.uiver.app/#q=WzAsNCxbMSwwLCJMXFxvZG90IE1cXG9kb3QgTlxcb2RvdCBTIl0sWzAsMCwiIFBcXG9kb3QgTCBcXG9kb3QgTVxcb2RvdCBOIl0sWzAsMSwiTFxcb2RvdCBRIFxcb2RvdCBNXFxvZG90IE4gIl0sWzEsMSwiTFxcb2RvdCBNXFxvZG90IFJcXG9kb3QgTiAiXSxbMSwyLCJ7XFxkZWx0YVxcb2RvdCBcXG1hdGhybXtpZH1fe01cXG9kb3QgTn19IiwyLHsibGV2ZWwiOjJ9XSxbMSwwLCJ7KEwsIFxcZGVsdGEpXFxib3hkb3QoTSwgXFxiZXRhKVxcYm94ZG90KE4sIFxcZ2FtbWEpfSIsMCx7ImxldmVsIjoyfV0sWzIsMywie1xcbWF0aHJte2lkfV9MXFxvZG90IFxcYmV0YVxcb2RvdCBcXG1hdGhybXtpZH1fTn0iLDIseyJsZXZlbCI6Mn1dLFsyLDAsIntcXG1hdGhybXtpZH1fTFxcb2RvdCAoKE0sIFxcYmV0YSlcXGJveGRvdChOLCBcXGdhbW1hKSl9IiwxLHsibGV2ZWwiOjJ9XSxbMywwLCJ7XFxtYXRocm17aWR9X0xcXG9kb3QgXFxtYXRocm17aWR9X01cXG9kb3QgXFxnYW1tYX0iLDIseyJsZXZlbCI6Mn1dXQ==
\adjustbox{scale=.9,center}{\begin{tikzcd}[column sep=1in,row sep=.5in]
	{ P\odot L \odot M\odot N} & {L\odot M\odot N\odot S} \\
	{L\odot Q \odot M\odot N } & {L\odot M\odot R\odot N }
	\arrow["{{(L, \delta)\boxdot(M, \beta)\boxdot(N, \gamma)}}", Rightarrow, from=1-1, to=1-2]
	\arrow["{{\delta\odot \mathrm{id}_{M\odot N}}}"', Rightarrow, from=1-1, to=2-1]
	\arrow["{{\mathrm{id}_L\odot ((M, \beta)\boxdot(N, \gamma))}}"{description}, Rightarrow, from=2-1, to=1-2]
	\arrow["{{\mathrm{id}_L\odot \beta\odot \mathrm{id}_N}}"', Rightarrow, from=2-1, to=2-2]
	\arrow["{{\mathrm{id}_L\odot \mathrm{id}_M\odot \gamma}}"', Rightarrow, from=2-2, to=1-2]
\end{tikzcd}}
\caption{$(L, \delta)\boxdot\left[(M, \beta)\boxdot(N, \gamma)\right]$}
\end{subfigure}
\caption{Compositions of three 1-cells in $\Lambda \B$}\label{eq:assoctivity_Lambda_B}
\end{figure}
\end{proof}

\begin{rmk}  In most of the following diagrams we will suppress structure morphisms as we did in \cref{eq:assoctivity_Lambda_B}.  This is justified by the coherence results in \cite{bicat_coherence,MP:coherence}.   We will retain them in examples like \eqref{eq:assocator_1-cells} where they are essential to understanding the diagram.
\end{rmk}

Let $P: V\to V$ be an object in $\Lambda \B$. Let $(U_V, r\ell^{-1}): P\to P$  be the 1-cell in $\Lambda \B$ equipped with the composition \[\xymatrix{P\odot U_V \ar@{=>}[r]^-r & P \ar@{=>}[r]^-{\ell^{-1}} & U_V \odot P,} \] where $U_V$ is the unit object for $V$ in $\B$ and $r, \ell$ are the right and left unitors in $\B$, respectively. 

\begin{lem}\label{LambdaB_unit} Let  $N: V\to R$ be a 1-cell in $\B$ and $(N, \gamma): P\to Q$ be a 1-cell in $\Lambda \B$.
The unitor maps 
\[r\colon N\odot U_R\to N\text{ and }\ell\colon U_V\odot  N\to N\] 
in $\B$ are 2-cells 
\[(N, \gamma)\boxdot (U_R, \ell^{-1}r) \Rightarrow (N, \gamma)\text{ and } (U_V, r^{-1}\ell)\boxdot (N, \gamma)\Rightarrow (N, \gamma)\] in $\Lambda \B$. 
\end{lem}
In particular, the unitors in $\B$ are unitors in $\Lambda\B$ satisfying the triangle axiom.


\begin{proof}
The commutative diagram \cref{existremunitor} shows the map $r$ is a 2-cell in $\Lambda \B$,
$\ltwo{r}: (N, \gamma)\boxdot (U_R, r\ell^{-1}) \Rightarrow (N, \gamma).$  The case of $\ell$ is similar.  
\end{proof}
\begin{figure}
% %https://q.uiver.app/#q=WzAsMTIsWzAsMCwiUFxcb2RvdCAoTlxcb2RvdCBVX1IpIl0sWzAsMSwiIChQXFxvZG90IE4pXFxvZG90IFVfUiJdLFsyLDAsIlBcXG9kb3QgTiJdLFswLDIsIiAoTlxcb2RvdCBRKVxcb2RvdCBVX1IiXSxbMiwxLCJQXFxvZG90IE4iXSxbMCwzLCIgTlxcb2RvdCAoUVxcb2RvdCBVX1IpIl0sWzIsMiwiTlxcb2RvdCBRIl0sWzAsNCwiIE5cXG9kb3QgKFVfUlxcb2RvdCBRKSJdLFswLDUsIiAoTlxcb2RvdCBVX1IpXFxvZG90IFEiXSxbMiw1LCJOXFxvZG90IFEiXSxbMiwzLCJOXFxvZG90IFEiXSxbMiw0LCJOXFxvZG90IFEiXSxbMSwwLCJhIiwwLHsibGV2ZWwiOjJ9XSxbMCwyLCJcXGlkXFxvZG90IHIiLDAseyJsZXZlbCI6Mn1dLFsxLDMsIlxcZ2FtbWFcXG9kb3RcXGlkIiwyLHsibGV2ZWwiOjJ9XSxbMSw0LCJyIiwwLHsibGV2ZWwiOjJ9XSxbMyw1LCJhIiwyLHsibGV2ZWwiOjJ9XSxbMyw2LCJyIiwwLHsibGV2ZWwiOjJ9XSxbNSw3LCJcXGlkXFxvZG90IHJcXGVsbF57LTF9IiwyLHsibGV2ZWwiOjJ9XSxbOCw3LCJhIiwwLHsibGV2ZWwiOjJ9XSxbOCw5LCJyXFxvZG90XFxpZCIsMCx7ImxldmVsIjoyfV0sWzIsNCwiIiwyLHsibGV2ZWwiOjIsInN0eWxlIjp7ImhlYWQiOnsibmFtZSI6Im5vbmUifX19XSxbNCw2LCJcXGdhbW1hIiwwLHsibGV2ZWwiOjJ9XSxbMCw0LCJcXHRleHR7TmF0dXJhbGl0eX0iLDEseyJzdHlsZSI6eyJib2R5Ijp7Im5hbWUiOiJub25lIn0sImhlYWQiOnsibmFtZSI6Im5vbmUifX19XSxbMyw0LCJcXHRleHR7TmF0LiBvZiBVbml0b3J9IiwxLHsic3R5bGUiOnsiYm9keSI6eyJuYW1lIjoibm9uZSJ9LCJoZWFkIjp7Im5hbWUiOiJub25lIn19fV0sWzYsMTAsIiIsMSx7ImxldmVsIjoyLCJzdHlsZSI6eyJoZWFkIjp7Im5hbWUiOiJub25lIn19fV0sWzUsMTAsIlxcaWRcXG9kb3QgciIsMCx7ImxldmVsIjoyfV0sWzUsNiwiXFx0ZXh0e1RyaWFuZ2xlIEF4aW9tfSIsMSx7InN0eWxlIjp7ImJvZHkiOnsibmFtZSI6Im5vbmUifSwiaGVhZCI6eyJuYW1lIjoibm9uZSJ9fX1dLFsxMSw5LCIiLDIseyJsZXZlbCI6Miwic3R5bGUiOnsiaGVhZCI6eyJuYW1lIjoibm9uZSJ9fX1dLFsxMCwxMSwiIiwyLHsibGV2ZWwiOjIsInN0eWxlIjp7ImhlYWQiOnsibmFtZSI6Im5vbmUifX19XSxbNywxMSwiXFxpZFxcb2RvdFxcZWxsIiwwLHsibGV2ZWwiOjJ9XSxbMTEsOCwiXFx0ZXh0e1RyaWFuZ2xlIEF4aW9tfSIsMSx7InN0eWxlIjp7ImJvZHkiOnsibmFtZSI6Im5vbmUifSwiaGVhZCI6eyJuYW1lIjoibm9uZSJ9fX1dLFs3LDEwLCI9IiwxLHsic3R5bGUiOnsiYm9keSI6eyJuYW1lIjoibm9uZSJ9fX1dXQ==&macro_url=https%3A%2F%2Fwww.overleaf.com%2Fproject%2F64c9c0fcea71c453e68a663b
\[\begin{tikzcd}
	{P\odot (N\odot U_R)} && {P\odot N} \\
	{ (P\odot N)\odot U_R} && {P\odot N} \\
	{ (N\odot Q)\odot U_R} && {N\odot Q} \\
	{ N\odot (Q\odot U_R)} && {N\odot Q} \\
	{ N\odot (U_R\odot Q)} && {N\odot Q} \\
	{ (N\odot U_R)\odot Q} && {N\odot Q}
	\arrow["{\id\odot r}", Rightarrow, from=1-1, to=1-3]
	\arrow["{\text{Naturality}}"{description}, draw=none, from=1-1, to=2-3]
	\arrow[equals, from=1-3, to=2-3]
	\arrow["a", Rightarrow, from=2-1, to=1-1]
	\arrow["r", Rightarrow, from=2-1, to=2-3]
	\arrow["{\gamma\odot\id}"', Rightarrow, from=2-1, to=3-1]
	\arrow["\gamma", Rightarrow, from=2-3, to=3-3]
	\arrow["{\text{Nat. of Unitor}}"{description}, draw=none, from=3-1, to=2-3]
	\arrow["r", Rightarrow, from=3-1, to=3-3]
	\arrow["a"', Rightarrow, from=3-1, to=4-1]
	\arrow[equals, from=3-3, to=4-3]
	\arrow["{\text{Triangle Axiom}}"{description}, draw=none, from=4-1, to=3-3]
	\arrow["{\id\odot r}", Rightarrow, from=4-1, to=4-3]
	\arrow["{\id\odot r\ell^{-1}}"', Rightarrow, from=4-1, to=5-1]
	\arrow[equals, from=4-3, to=5-3]
	\arrow["{=}"{description}, no body, from=5-1, to=4-3]
	\arrow["{\id\odot\ell}", Rightarrow, from=5-1, to=5-3]
	\arrow["{\text{Triangle Axiom}}"{description}, draw=none, from=5-3, to=6-1]
	\arrow[equals, from=5-3, to=6-3]
	\arrow["a", Rightarrow, from=6-1, to=5-1]
	\arrow["{r\odot\id}", Rightarrow, from=6-1, to=6-3]
\end{tikzcd}\]
\caption{Existence of unitor for $\Lambda \B$}\label{existremunitor}
\end{figure}

 \cref{LambdaBPQ_Cat,LambdaB_assoc,LambdaB_unit} complete the proof of \cref{Lambda_bicat} that $\Lambda\B$ is a bicategory. 

\section{Functors induced on twisted endomorphism bicategories}\label{sec:remnant_functors} 
In this section we verify that the construction of \cref{sec:remnant} extends to functors. 

\begin{lem}\label{inducted_functor_lem_1} The 2-cell $\phi \colon F(M)\odot F(N)\to F(M\odot N)$ in $\B$ is a 2-cell
\[\Lambda F(M,\beta)\boxdot \Lambda F(N,\gamma)\to \Lambda F((M,\beta)\boxdot (N,\gamma))\]
 in $\Lambda \B$.

\end{lem}


\begin{proof}

\begin{figure}
% https://q.uiver.app/#q=WzAsMTIsWzAsMCwiRihQKVxcb2RvdCBGKE0pXFxvZG90IEYoTikiXSxbMSwxLCJGKFBcXG9kb3QgTSlcXG9kb3QgRihOKSJdLFsxLDIsIkYoTVxcb2RvdCBRKVxcb2RvdCBGKE4pIl0sWzEsMywiRihNKVxcb2RvdCBGKFEpXFxvZG90IEYoTikiXSxbMSw0LCJGKE0pXFxvZG90IEYoUVxcb2RvdCBOKSJdLFswLDYsIkYoTSlcXG9kb3QgRihOKVxcb2RvdCBGKFQpIl0sWzEsNSwiRihNKVxcb2RvdCBGKE5cXG9kb3QgVCkiXSxbMiwxLCJGKFBcXG9kb3QgTVxcb2RvdCBOKSJdLFsyLDMsIkYoTVxcb2RvdCBRXFxvZG90IE4pIl0sWzIsNSwiRihNXFxvZG90IE5cXG9kb3QgVCkiXSxbMywwLCJGKFApXFxvZG90IEYoTVxcb2RvdCBOKSJdLFszLDYsIkYoTVxcb2RvdCBOKVxcb2RvdCBGKFQpIl0sWzAsMSwiXFxwaGlcXG9kb3QgXFxtYXRocm17aWR9IiwyLHsibGV2ZWwiOjJ9XSxbMCw1LCJcXExhbWJkYSBGKE0sXFxiZXRhKVxcYm94ZG90IFxcTGFtYmRhIEYoTixcXGdhbW1hKSIsMix7ImxldmVsIjoyLCJzdHlsZSI6eyJib2R5Ijp7Im5hbWUiOiJkb3R0ZWQifX19XSxbMSwyLCJGKFxcYmV0YSlcXG9kb3QgXFxtYXRocm17aWR9IiwwLHsibGV2ZWwiOjJ9XSxbMywyLCJcXHBoaVxcb2RvdCBcXG1hdGhybXtpZH0iLDIseyJsZXZlbCI6Mn1dLFszLDQsIlxcbWF0aHJte2lkfVxcb2RvdCBcXHBoaSIsMCx7ImxldmVsIjoyfV0sWzQsNiwiXFxtYXRocm17aWR9XFxvZG90IEYoXFxnYW1tYSkiLDAseyJsZXZlbCI6Mn1dLFs1LDYsIlxcbWF0aHJte2lkfVxcb2RvdCBcXHBoaSIsMCx7ImxldmVsIjoyfV0sWzcsOCwiRihcXGJldGFcXG9kb3QgXFxtYXRocm17aWR9KSIsMCx7ImxldmVsIjoyfV0sWzgsOSwiRihcXG1hdGhybXtpZH1cXG9kb3QgXFxnYW1tYSkiLDAseyJsZXZlbCI6Mn1dLFs2LDksIlxccGhpIiwyLHsibGV2ZWwiOjJ9XSxbMSw3LCJcXHBoaSIsMCx7ImxldmVsIjoyfV0sWzExLDksIlxccGhpIiwyLHsibGV2ZWwiOjJ9XSxbMTAsNywiXFxwaGkiLDAseyJsZXZlbCI6Mn1dLFswLDEwLCJcXGlkIFxcb2RvdCBcXHBoaSIsMCx7ImxldmVsIjoyfV0sWzEwLDExLCJcXExhbWJkYSBGKChNLFxcYmV0YSlcXGJveGRvdChOLFxcZ2FtbWEpKSIsMCx7ImxldmVsIjoyLCJzdHlsZSI6eyJib2R5Ijp7Im5hbWUiOiJkb3R0ZWQifX19XSxbNSwxMSwiXFxwaGlcXG9kb3QgXFxpZCIsMix7ImxldmVsIjoyfV0sWzIsOCwiXFxwaGkiLDAseyJsZXZlbCI6Mn1dLFs0LDgsIlxccGhpIiwyLHsibGV2ZWwiOjJ9XV0=
\[\begin{tikzcd}[column sep= .0in, row sep =. 35in]
	{F(P)\odot F(M)\odot F(N)} &&& {F(P)\odot F(M\odot N)} \\
	& {F(P\odot M)\odot F(N)} & {F(P\odot M\odot N)} \\
	& {F(M\odot Q)\odot F(N)} \\
	& {F(M)\odot F(Q)\odot F(N)} & {F(M\odot Q\odot N)} \\
	& {F(M)\odot F(Q\odot N)} \\
	& {F(M)\odot F(N\odot T)} & {F(M\odot N\odot T)} \\
	{F(M)\odot F(N)\odot F(T)} &&& {F(M\odot N)\odot F(T)}
	\arrow["{\id \odot \phi}", Rightarrow,  from=1-1, to=1-4]
	\arrow["{\phi\odot \mathrm{id}}"', Rightarrow,  from=1-1, to=2-2]
	\arrow["{\Lambda F(M,\beta)\boxdot \Lambda F(N,\gamma)}"', Rightarrow,  dotted, from=1-1, to=7-1]
	\arrow["\phi", Rightarrow,  from=1-4, to=2-3]
	\arrow["{\Lambda F((M,\beta)\boxdot(N,\gamma))}", Rightarrow,  dotted, from=1-4, to=7-4]
	\arrow["\phi", Rightarrow,  from=2-2, to=2-3]
	\arrow["{F(\beta)\odot \mathrm{id}}", Rightarrow,  from=2-2, to=3-2]
	\arrow["{F(\beta\odot \mathrm{id})}", Rightarrow,  from=2-3, to=4-3]
	\arrow["\phi", Rightarrow,  from=3-2, to=4-3]
	\arrow["{\phi\odot \mathrm{id}}"', Rightarrow,  from=4-2, to=3-2]
	\arrow["{\mathrm{id}\odot \phi}", Rightarrow,  from=4-2, to=5-2]
	\arrow["{F(\mathrm{id}\odot \gamma)}", Rightarrow,  from=4-3, to=6-3]
	\arrow["\phi"', Rightarrow,  from=5-2, to=4-3]
	\arrow["{\mathrm{id}\odot F(\gamma)}", Rightarrow,  from=5-2, to=6-2]
	\arrow["\phi"', Rightarrow,  from=6-2, to=6-3]
	\arrow["{\mathrm{id}\odot \phi}", Rightarrow,  from=7-1, to=6-2]
	\arrow["{\phi\odot \id}"', Rightarrow,  from=7-1, to=7-4]
	\arrow["\phi"', Rightarrow,  from=7-4, to=6-3]
\end{tikzcd}\]
\caption{Compatibility of $\Lambda F$ and $\boxdot$}\label{Lambda_F_boxdot}
\end{figure}
We show that $\phi \colon F(M)\odot F(N)\to F(M\odot N)$ defines a 2-cell in $\Lambda \B$ by showing that the diagram in \eqref{boxdot_Lambda} commutes.
\begin{equation}\label{boxdot_Lambda}
% https://q.uiver.app/#q=WzAsNCxbMCwxLCJGKFApXFxvZG90IEYoTVxcb2RvdCBOKSJdLFszLDEsIkYoTVxcb2RvdCBOKVxcb2RvdCBGKFQpIl0sWzAsMCwiRihQKVxcb2RvdCAoRihNKVxcb2RvdCBGKE4pKSJdLFszLDAsIihGKE0pXFxvZG90IEYoTikpXFxvZG90IEYoVCkiXSxbMCwxLCJcXExhbWJkYSBGKChNLFxcYmV0YSlcXGJveGRvdCAoTixcXGdhbW1hKSkiLDAseyJsZXZlbCI6Miwic3R5bGUiOnsiYm9keSI6eyJuYW1lIjoiZG90dGVkIn19fV0sWzIsMywiXFxMYW1iZGEgRihNLFxcYmV0YSlcXGJveGRvdCBcXExhbWJkYSBGKE4sXFxnYW1tYSkiLDAseyJsZXZlbCI6Miwic3R5bGUiOnsiYm9keSI6eyJuYW1lIjoiZG90dGVkIn19fV0sWzIsMCwiXFxpZFxcb2RvdFxccGhpIiwyLHsibGV2ZWwiOjJ9XSxbMywxLCJcXHBoaVxcb2RvdFxcaWQiLDAseyJsZXZlbCI6Mn1dXQ==&macro_url=https%3A%2F%2Fwww.overleaf.com%2Fproject%2F64c9c0fcea71c453e68a663b
\begin{tikzcd}
	{F(P)\odot (F(M)\odot F(N))} &&& {(F(M)\odot F(N))\odot F(T)} \\
	{F(P)\odot F(M\odot N)} &&& {F(M\odot N)\odot F(T)}
	\arrow["{\Lambda F(M,\beta)\boxdot \Lambda F(N,\gamma)}", Rightarrow, dotted, from=1-1, to=1-4]
	\arrow["{\id\odot\phi}"', Rightarrow, from=1-1, to=2-1]
	\arrow["{\phi\odot\id}", Rightarrow, from=1-4, to=2-4]
	\arrow["{\Lambda F((M,\beta)\boxdot (N,\gamma))}", Rightarrow, dotted, from=2-1, to=2-4]
\end{tikzcd}
\end{equation}
See \cref{Lambda_F_boxdot}.
\end{proof}

\begin{lem}\label{inducted_functor_lem_2}
 For a 0-cell $V$ in $\Lambda \B$, the 2-cell $\iota\colon U_{F(V)}\to F(U_V)$ is a 2-cell $(U_{F(V)},r\ell^{-1})\to \Lambda F(U_V,r\ell^{-1})$.
\end{lem}

\begin{proof}
The image of $(U_V,r\ell^{-1})$ under $\Lambda F$ is the pair 
\[(F(U_V), F(P)\odot F(U_V)\Rightarrow F(P\odot U_V)
\Rightarrow F(P) \Rightarrow F(U_V\odot P)\Rightarrow F(U_V)\odot F(P)).\]
Since the diagram in \cref{lambda_F_unit} commutes, 
$\iota$ defines a map in $\Lambda \cc$.
\begin{figure}
\centering
% https://q.uiver.app/#q=WzAsOCxbMSwwLCJGKFApXFxvZG90IEYoVV9WKSJdLFsxLDEsIkYoUFxcb2RvdCBVX1YpIl0sWzEsMiwiRihQKSJdLFsxLDMsIkYoVV9WXFxvZG90IFApIl0sWzEsNCwiRihVX1YpXFxvZG90IEYoUCkiXSxbMCwwLCJGKFApXFxvZG90IFVfe0YoVil9Il0sWzAsNCwiVV97RihWKX1cXG9kb3QgRihQKSJdLFswLDIsIkYoUCkiXSxbNywyLCIiLDIseyJsZXZlbCI6Miwic3R5bGUiOnsiaGVhZCI6eyJuYW1lIjoibm9uZSJ9fX1dLFs1LDcsInIiLDIseyJsZXZlbCI6Mn1dLFs3LDYsIlxcZWxsXnstMX0iLDIseyJsZXZlbCI6Mn1dLFs1LDAsIlxcbWF0aHJte2lkfVxcb2RvdCBcXGlvdGEiLDIseyJsZXZlbCI6Mn1dLFswLDEsIlxccGhpIiwwLHsibGV2ZWwiOjJ9XSxbMSwyLCJGKHIpIiwwLHsibGV2ZWwiOjJ9XSxbMiwzLCJGKFxcZWxsXnstMX0pIiwwLHsibGV2ZWwiOjJ9XSxbMyw0LCJcXHBoaV57LTF9IiwwLHsibGV2ZWwiOjJ9XSxbNiw0LCJcXGlvdGFcXG9kb3QgXFxtYXRocm17aWR9IiwyLHsibGV2ZWwiOjJ9XV0=
\begin{tikzcd}
	{F(P)\odot U_{F(V)}} & {F(P)\odot F(U_V)} \\
	& {F(P\odot U_V)} \\
	{F(P)} & {F(P)} \\
	& {F(U_V\odot P)} \\
	{U_{F(V)}\odot F(P)} & {F(U_V)\odot F(P)}
	\arrow[Rightarrow, no head, from=3-1, to=3-2]
	\arrow["r"', Rightarrow, from=1-1, to=3-1]
	\arrow["{\ell^{-1}}"', Rightarrow, from=3-1, to=5-1]
	\arrow["{\mathrm{id}\odot \iota}"', Rightarrow, from=1-1, to=1-2]
	\arrow["\phi", Rightarrow, from=1-2, to=2-2]
	\arrow["{F(r)}", Rightarrow, from=2-2, to=3-2]
	\arrow["{F(\ell^{-1})}", Rightarrow, from=3-2, to=4-2]
	\arrow["{\phi^{-1}}", Rightarrow, from=4-2, to=5-2]
	\arrow["{\iota\odot \mathrm{id}}"', Rightarrow, from=5-1, to=5-2]
\end{tikzcd}
\caption{$\Lambda F$ respects the unit}\label{lambda_F_unit}
\end{figure}
\end{proof}

\lmabdaofacomposite*

\begin{proof}
   If $F$ is a strong functor, $F(M)$ is dualizable.   Then the pair 
    \begin{equation}\label{1_cell_Lambda_F}(F(M), \phi\circ F(\psi)\circ \phi^{-1}).
    \end{equation} 
 is a 1-cell in $\Lambda \mathcal{C}$. 
By \cref{inducted_functor_lem_1,inducted_functor_lem_2},
the coherence diagrams follow from the coherence diagrams for $F$.

    The compatibility with composition is straightforward to check. 
    On 0- and 2-cells, $\Lambda F=F$ and there is nothing to prove.  For 1-cells, 
    \begin{align*}
        (\Lambda (F\circ G))(M,\beta)
        &=((F\circ G)(M), \phi_{F\circ G} \circ ((F\circ G)(\beta)) \circ \phi_{F\circ G}^{-1})\\
        &=((F\circ G)(M), \phi_{F}\circ F(\phi_{G})\circ ((F\circ G)(\beta)) \circ F(\phi_{G}^{-1})\circ \phi_F^{-1})\\
        &=(\Lambda F)(G(M), \phi_{G}\circ G(\beta) \circ \phi_{G}^{-1})
        \\
        &=(\Lambda F)(\Lambda G)(M,\beta)
    \end{align*}
\end{proof}



\section{Traces on twisted endomorphism bicategories and characters}\label{sec:traces}
In this section we show  the bicategorical trace can be upgraded to a functor on the endomorphism bicategory of \cref{sec:remnant} and use this functor to define characters for 2-representations.
\thmcharacter*



To prove this theorem we first note the following consequence of  \cite[Proposition 7.5]{Ponto_2012}.

\begin{lem} \label{lem:comp}
    For the functor  $\rchi$ defined above and composable 1-cells $(M, \beta)$ and $(N, \gamma)$ in $\Lambda \B$, 
    \[\rchi[(N, \gamma)\boxdot (M, \beta)] = \rchi[(N, \gamma)]\circ \rchi[(M, \beta)].\] 
\end{lem}

We now verify that $\chi$ preserves unit 1-cells.

 \begin{lem} \label{lem:unit}
     For 
     the unit 1-cell $(U_V, r\ell^{-1}): P\to P$ in $\Lambda \B$, 
     \[\rchi[(U_V, r\ell^{-1})] = \id_{\sh{P}}.\]
 \end{lem}

 \begin{proof}
Recall that $(U_R, U_R)$ forms a dual pair, with coevaluation 
$\eta = \ell^{-1}=r^{-1}$ and evaluation 
$\epsilon = \ell=r$. 

In the diagram in \cref{fig:lambda_functor_unit}, $\rchi[(U_V, r\ell^{-1})]$ is the  top, right, and bottom composite.  
% https://q.uiver.app/#q=WzAsMTMsWzAsMCwiXFxzaHtQfSJdLFsxLDAsIlxcc2h7UFxcb2RvdCBVX1J9Il0sWzMsMCwiXFxzaHtQXFxvZG90KFVfUlxcb2RvdCBVX1IpfSJdLFszLDIsIlxcc2h7UFxcb2RvdCBVX1J9Il0sWzMsMywiXFxzaHsoVV9SXFxvZG90IFApXFxvZG90IFVfUn0iXSxbMyw0LCJcXHNoe1VfUlxcb2RvdCAoVV9SXFxvZG90IFApfSJdLFszLDUsIlxcc2h7KFVfUlxcb2RvdCBVX1IpXFxvZG90IFB9Il0sWzAsNSwiXFxzaHtQfSJdLFszLDEsIlxcc2h7KFBcXG9kb3QgVV9SKVxcb2RvdCBVX1J9Il0sWzEsNSwiXFxzaHtVX1JcXG9kb3QgUH0iXSxbMSwyLCJcXHNoe1BcXG9kb3QgVV9SfSJdLFsxLDMsIlxcc2h7UFxcb2RvdCBVX1J9Il0sWzEsNCwiXFxzaHtVX1JcXG9kb3QgUH0iXSxbMCwxLCJcXHNoe3Jeey0xfX0iLDAseyJsZXZlbCI6Mn1dLFsxLDIsIlxcc2h7XFxpZFxcb2RvdCBcXGV0YX0iLDAseyJjdXJ2ZSI6LTQsImxldmVsIjoyfV0sWzMsNCwiXFxzaHtcXGVsbF57LTF9XFxvZG90XFxpZH0iLDAseyJsZXZlbCI6Mn1dLFs0LDUsIlxcdGhldGEiLDAseyJsZXZlbCI6Mn1dLFs2LDUsIlxcc2h7YV57LTF9fSIsMix7ImxldmVsIjoyfV0sWzgsMiwiXFxzaHthXnstMX19IiwyLHsibGV2ZWwiOjJ9XSxbOCwzLCJcXHNoe3JcXG9kb3RcXGlkfSIsMCx7ImxldmVsIjoyfV0sWzYsOSwiXFxzaHtcXGVwc2lsb25cXG9kb3RcXGlkfSIsMCx7ImN1cnZlIjotNCwibGV2ZWwiOjJ9XSxbOSw3LCJcXHNoe1xcZWxsfSIsMix7ImxldmVsIjoyfV0sWzEsMiwiXFxzaHtcXGlkXFxvZG90IFxcZWxsXnstMX19IiwyLHsibGV2ZWwiOjJ9XSxbMCw3LCJcXGlkIiwwLHsibGV2ZWwiOjIsInN0eWxlIjp7ImhlYWQiOnsibmFtZSI6Im5vbmUifX19XSxbNiw5LCJcXHNoe3JcXG9kb3RcXGlkfSIsMix7ImxldmVsIjoyfV0sWzEsMTAsIlxcaWQiLDIseyJsZXZlbCI6Miwic3R5bGUiOnsiaGVhZCI6eyJuYW1lIjoibm9uZSJ9fX1dLFsxMCwzLCJcXGlkIiwyLHsibGV2ZWwiOjIsInN0eWxlIjp7ImhlYWQiOnsibmFtZSI6Im5vbmUifX19XSxbMTAsMTEsIlxcaWQiLDAseyJsZXZlbCI6Miwic3R5bGUiOnsiaGVhZCI6eyJuYW1lIjoibm9uZSJ9fX1dLFs0LDExLCJcXHNoe1xcZWxsXFxvZG90IFxcaWR9IiwyLHsibGV2ZWwiOjJ9XSxbMTEsMTIsIlxcdGhldGEiLDIseyJsZXZlbCI6Mn1dLFs1LDEyLCJcXHNoe1xcaWRcXG9kb3QgXFxlbGx9IiwyLHsibGV2ZWwiOjJ9XSxbMTIsOSwiXFxpZCIsMix7ImxldmVsIjoyLCJzdHlsZSI6eyJoZWFkIjp7Im5hbWUiOiJub25lIn19fV1d&macro_url=https%3A%2F%2Fwww.overleaf.com%2Fproject%2F64c9c0fcea71c453e68a663b
\begin{figure}
\begin{tikzcd}
	{\sh{P}} & {\sh{P\odot U_R}} && {\sh{P\odot(U_R\odot U_R)}} \\
	&&& {\sh{(P\odot U_R)\odot U_R}} \\
	& {\sh{P\odot U_R}} && {\sh{P\odot U_R}} \\
	& {\sh{P\odot U_R}} && {\sh{(U_R\odot P)\odot U_R}} \\
	& {\sh{U_R\odot P}} && {\sh{U_R\odot (U_R\odot P)}} \\
	{\sh{P}} & {\sh{U_R\odot P}} && {\sh{(U_R\odot U_R)\odot P}}
	\arrow["{\sh{r^{-1}}}", Rightarrow, from=1-1, to=1-2]
	\arrow["\id", equals, from=1-1, to=6-1]
	\arrow["{\sh{\id\odot \eta}}", curve={height=-24pt}, Rightarrow, from=1-2, to=1-4]
	\arrow["{\sh{\id\odot \ell^{-1}}}"', Rightarrow, from=1-2, to=1-4]
	\arrow["\id"', equals, from=1-2, to=3-2]
	\arrow["{\sh{a^{-1}}}"', Rightarrow, from=2-4, to=1-4]
	\arrow["{\sh{r\odot\id}}", Rightarrow, from=2-4, to=3-4]
	\arrow["\id"', equals, from=3-2, to=3-4]
	\arrow["\id", equals, from=3-2, to=4-2]
	\arrow["{\sh{\ell^{-1}\odot\id}}", Rightarrow, from=3-4, to=4-4]
	\arrow["\theta"', Rightarrow, from=4-2, to=5-2]
	\arrow["{\sh{\ell\odot \id}}"', Rightarrow, from=4-4, to=4-2]
	\arrow["\theta", Rightarrow, from=4-4, to=5-4]
	\arrow["\id"', equals, from=5-2, to=6-2]
	\arrow["{\sh{\id\odot \ell}}"', Rightarrow, from=5-4, to=5-2]
	\arrow["{\sh{\ell}}"', Rightarrow, from=6-2, to=6-1]
	\arrow["{\sh{a^{-1}}}"', Rightarrow, from=6-4, to=5-4]
	\arrow["{\sh{\epsilon\odot\id}}", curve={height=-24pt}, Rightarrow, from=6-4, to=6-2]
	\arrow["{\sh{r\odot\id}}"', Rightarrow, from=6-4, to=6-2]
\end{tikzcd}
\caption{ $\rchi[(U_V, r\ell^{-1})]$ }\label{fig:lambda_functor_unit}
\end{figure}
The regions of the diagram commute by coherence for bicategories and shadows \cite{MP:coherence}.
 \end{proof}

To consider the behavior on 2-cells we need a preliminary lemma. 
\begin{lem} \label{lem:treq}
    Let $\B$ be a bicategory with shadow functors $\sh{\text{-}}$. 
    Given a dualizable 1-cell $M$, isomorphism 2-cell $\phi\colon M\Rightarrow N$, and commuting diagram 
    % https://q.uiver.app/#q=WzAsNCxbMCwwLCJQXFxvZG90IE0iXSxbMSwwLCJNXFxvZG90IFEiXSxbMCwxLCJQXFxvZG90IE4iXSxbMSwxLCJOXFxvZG90IFEiXSxbMCwyLCJcXGlkIFxcb2RvdCBcXHBoaSIsMix7ImxldmVsIjoyfV0sWzAsMSwiXFxiZXRhIiwwLHsibGV2ZWwiOjJ9XSxbMiwzLCJcXGdhbW1hIiwwLHsibGV2ZWwiOjJ9XSxbMSwzLCJcXHBoaVxcb2RvdCBcXGlkIiwwLHsibGV2ZWwiOjJ9XV0=&macro_url=https%3A%2F%2Fwww.overleaf.com%2Fproject%2F64c9c0fcea71c453e68a663b
\[\begin{tikzcd}
	{P\odot M} & {M\odot Q} \\
	{P\odot N} & {N\odot Q}
	\arrow["\beta", Rightarrow, from=1-1, to=1-2]
	\arrow["{\id \odot \phi}"', Rightarrow, from=1-1, to=2-1]
	\arrow["{\phi\odot \id}", Rightarrow, from=1-2, to=2-2]
	\arrow["\gamma", Rightarrow, from=2-1, to=2-2]
\end{tikzcd}\]
     then 
   $\tr(\beta) = \tr(\gamma).$
\end{lem}

\begin{proof}
       Let  $\eta_M$ and $\epsilon_M$ be a choice of coevaluation and evaluation for $M$ with dual $M^\ast$.
    Since $N$ is isomorphic to $M$ and $M$ is dualizable, $N$ is dualizable and a choice of coevaluation and evaluation for $N$ are 
    \[\xymatrix{\eta: U_R \ar@{=>}[r]^-{\eta_M} & M\odot M^\ast \ar@{=>}[r]^{\phi\odot\id} & N\odot M^\ast}\] and
    \[\xymatrix{\epsilon: M^\ast \odot N \ar@{=>}[r]^-{\id\odot\phi^{-1}} & M^\ast\odot M \ar@{=>}[r]^-{\epsilon_M} & U_S}.\]


    Since traces are independent of the choice of dual \cite[4.5.2]{p:thesis}, we can use this dual pair to compute the trace of $\gamma$. In the commuting diagram in \cref{comparison_of_traces}, the left composite is the trace of $\beta$ and the right composite is the trace of $\gamma$.
%
\begin{figure}
\centering
% https://q.uiver.app/#q=WzAsMTIsWzAsMCwiXFxzaHtQfSJdLFswLDEsIlxcc2h7UFxcb2RvdCBVX1N9ICJdLFsxLDEsIlxcc2h7UFxcb2RvdCBVX1N9ICJdLFswLDIsIiBcXHNoe1BcXG9kb3QgTVxcb2RvdCBNXlxcYXN0fSJdLFswLDMsIlxcc2h7TVxcb2RvdCBRXFxvZG90IE1eXFxhc3R9Il0sWzIsMiwiXFxzaHtQXFxvZG90IE5cXG9kb3QgTV5cXGFzdH0iXSxbMCw0LCJcXHNoe1FcXG9kb3QgTV5cXGFzdFxcb2RvdCBNfSAiXSxbMiwzLCJcXHNoe05cXG9kb3QgUVxcb2RvdCBNXlxcYXN0fSJdLFswLDUsIlxcc2h7USBcXG9kb3QgVV9TfSJdLFsyLDQsIlxcc2h7UVxcb2RvdCBNXlxcYXN0XFxvZG90IE59Il0sWzAsNiwiXFxzaHtRfSJdLFsxLDUsIlxcc2h7USBcXG9kb3QgVV9TfSJdLFswLDEsIlxcc2h7cl57LTF9fSIsMix7ImxldmVsIjoyfV0sWzAsMiwiXFxzaHtyXnstMX19IiwwLHsibGV2ZWwiOjJ9XSxbMSwzLCJcXHNoe1xcaWRcXG9kb3RcXGV0YV9NfSIsMix7ImxldmVsIjoyfV0sWzEsMiwiIiwwLHsibGV2ZWwiOjIsInN0eWxlIjp7ImhlYWQiOnsibmFtZSI6Im5vbmUifX19XSxbMyw0LCJcXHNoe1xcYmV0YVxcb2RvdFxcaWR9IiwyLHsibGV2ZWwiOjJ9XSxbMyw1LCJcXHNoe1xcaWRcXG9kb3RcXHBoaVxcb2RvdFxcaWR9IiwwLHsibGV2ZWwiOjJ9XSxbNCw2LCJcXHRoZXRhIiwyLHsibGV2ZWwiOjJ9XSxbNCw3LCJcXHNoe1xccGhpXFxvZG90XFxpZFxcb2RvdFxcaWR9IiwwLHsibGV2ZWwiOjJ9XSxbNiw4LCJcXHNoe1xcaWRcXG9kb3RcXGVwc2lsb25fTX0iLDIseyJsZXZlbCI6Mn1dLFs2LDksIlxcc2h7XFxpZFxcb2RvdFxcaWRcXG9kb3RcXHBoaX0iLDAseyJsZXZlbCI6Mn1dLFs4LDEwLCJcXHNoe3J9IiwyLHsibGV2ZWwiOjJ9XSxbOCwxMSwiIiwwLHsibGV2ZWwiOjIsInN0eWxlIjp7ImhlYWQiOnsibmFtZSI6Im5vbmUifX19XSxbMiw1LCJcXHNoe1xcaWRcXG9kb3RcXGV0YX0iLDAseyJsZXZlbCI6Miwic3R5bGUiOnsiYm9keSI6eyJuYW1lIjoiZG90dGVkIn19fV0sWzUsNywiXFxzaHtcXGdhbW1hXFxvZG90XFxpZH0iLDAseyJsZXZlbCI6Mn1dLFs3LDksIlxcdGhldGEiLDAseyJsZXZlbCI6Mn1dLFs5LDExLCJcXHNoe1xcaWRcXG9kb3RcXGVwc2lsb259IiwwLHsibGV2ZWwiOjIsInN0eWxlIjp7ImJvZHkiOnsibmFtZSI6ImRvdHRlZCJ9fX1dLFsxMSwxMCwiXFxzaHtyfSIsMCx7ImxldmVsIjoyfV1d
\begin{tikzcd}
	{\sh{P}} \\
	{\sh{P\odot U_S} } & {\sh{P\odot U_S} } \\
	{ \sh{P\odot M\odot M^\ast}} && {\sh{P\odot N\odot M^\ast}} \\
	{\sh{M\odot Q\odot M^\ast}} && {\sh{N\odot Q\odot M^\ast}} \\
	{\sh{Q\odot M^\ast\odot M} } && {\sh{Q\odot M^\ast\odot N}} \\
	{\sh{Q \odot U_S}} & {\sh{Q \odot U_S}} \\
	{\sh{Q}}
	\arrow["{\sh{r^{-1}}}"', Rightarrow, from=1-1, to=2-1]
	\arrow["{\sh{r^{-1}}}", Rightarrow, from=1-1, to=2-2]
	\arrow["{\sh{\id\odot\eta_M}}"', Rightarrow, from=2-1, to=3-1]
	\arrow[Rightarrow, no head, from=2-1, to=2-2]
	\arrow["{\sh{\beta\odot\id}}"', Rightarrow, from=3-1, to=4-1]
	\arrow["{\sh{\id\odot\phi\odot\id}}", Rightarrow, from=3-1, to=3-3]
	\arrow["\theta"', Rightarrow, from=4-1, to=5-1]
	\arrow["{\sh{\phi\odot\id\odot\id}}", Rightarrow, from=4-1, to=4-3]
	\arrow["{\sh{\id\odot\epsilon_M}}"', Rightarrow, from=5-1, to=6-1]
	\arrow["{\sh{\id\odot\id\odot\phi}}", Rightarrow, from=5-1, to=5-3]
	\arrow["{\sh{r}}"', Rightarrow, from=6-1, to=7-1]
	\arrow[Rightarrow, no head, from=6-1, to=6-2]
	\arrow["{\sh{\id\odot\eta}}", Rightarrow, dotted, from=2-2, to=3-3]
	\arrow["{\sh{\gamma\odot\id}}", Rightarrow, from=3-3, to=4-3]
	\arrow["\theta", Rightarrow, from=4-3, to=5-3]
	\arrow["{\sh{\id\odot\epsilon}}", Rightarrow, dotted, from=5-3, to=6-2]
	\arrow["{\sh{r}}", Rightarrow, from=6-2, to=7-1]
\end{tikzcd}
\caption{Comparison of traces (\cref{lem:treq})}\label{comparison_of_traces}
\end{figure}
\end{proof}

Composition of 2-cells is preserved by the map $\rchi$ follows from the limited number of 2-cells in the target category.

 \begin{lem} \label{lem:2comp}
     For the map $\rchi$ defined above and composable 2-cells $\ltwo{\phi}: (M, \beta) \Rightarrow (N, \gamma)$ and $\ltwo{\psi}: (N, \gamma) \Rightarrow (L, \sigma)$, 
     \[\rchi(\ltwo{\phi}\Box \ltwo{\psi}) = \rchi(\ltwo{\phi})\circ \rchi(\ltwo{\psi}).\]
 \end{lem}

\begin{proof}
    $\rchi(\ltwo{\phi})$ is a 2-cell  $\tr(\beta)\Rightarrow \tr(\gamma)$. The only 2-cells in $T$ are identities. By \cref{lem:treq} and $M\cong N\cong L$, $\tr(\beta)=\tr(\gamma)=\tr(\sigma)$, which means $\id_{\tr(\beta)} = \id_{\tr(\gamma)} = \id_{\tr(\sigma)}$. 
\end{proof}


\noindent
\begin{proof}[Proof of \cref{thm:character}]
    \cref{lem:comp} shows that $\rchi$ respects 1-cell composition. \cref{lem:unit} shows that $\rchi$ preserves the unit 1-cell. \cref{lem:2comp} shows that $\rchi$ respects 2-cell composition. $\rchi$ maps 2-cells in $\Lambda\B$ to identities, so $\rchi$ preserves the unit 2-cell trivially. 
\end{proof}
We now use the functor in \cref{thm:character} to define the character of a strong 2-functor (and so a 2-representation).

\defcharacter*


Conjugation invariance is fundamental to this definition.  In the case that $\mathcal{C}=BG$, the 1-cells in $\Lambda BG$ are equalities \[gh=hk\] for elements $g$, $h$, and $k$ of $G$, and 
traces associated to these equalities are the fundamental objects of interest.  For other choices of $\mathcal{C}$, the morphisms of $\Lambda \mathcal{C}$ are 2-cells 
\[P\odot M\to M\odot Q\] which we consider to be generalized conjugation relations between 1-cells. 

\section{Coshadows, cotraces, and cotwisted endomorphism bicategories}\label{sec:coshadows}
This section has two goals.  The first is to show that the approach above to characters in terms of the bicategorical trace extends to the bicategorical cotrace \cite{barhite2023bicategorical}.  The second goal, and the reason for the relevance of the first, is to more clearly illuminate how the approach of \cite{GK} sits inside the approach here.

Fundamental to the Bartlett and  Ganter--Kapranov character is the following definition.
\begin{defn}\cite[4.4]{Bartlett}\cite[3.1]{GK} Let $F\colon x\to x$ be an endo-1-cell in a bicategory $\mathcal{B}$. Then the {\bf categorical trace} of $F$ is \[\Tr(F)= 2\Hom_{\mathcal{B}}(U_x, F).\] 
\end{defn}
That is, the categorical trace of $F$ is the set of 2-cells from the identity 1-cell on $x$ to $F$.   

\begin{ex} The case of $2\Vect_k$ is particularly illuminating.  For a 1-cell $A = [A_{ij}]$ in $2\Vect_k$, 
\[\Tr(A) = 2\Hom_{2\Vect_k}(1_{[n]}, A) \cong \overset{n}{\underset{i=1}{\bigoplus}}\ A_{ii}.\]
\end{ex}

The categorical trace is not a shadow in the sense of \cref{defn:shadow}, but it a coshadow as defined by Barhite in \cite{barhite2023bicategorical}.   (See \cite[Example 4.4]{barhite2023bicategorical}.)
In this section we briefly describe the generalization of the approach above to characters using shadows to coshadows.  The arguments are the same, but this gives an approach in closer alignment to that in \cite{GK}.

Coshadows have the same properties as  shadows but with respect to the internal hom rather than the bicategorical composition.  A closed 
$\B$ bicategory  is equipped with left and right internal hom-functors
\[
	- \triangleleft - : \B(R, T) \times \B(R, S)^{op} \to \B(S, T)
\]
and
\[
	- \triangleright - : \mathscr{B}(S, T)^{op} \times \mathscr{B}(R, T) \to \mathscr{B}(R, S)
\]
for all triples of 0-cells $R, S, T$ and natural isomorphisms
\begin{equation}
\label{eq:closed_bicat_tensor_hom}
	\mathscr{B}(S, T)(N, P \triangleleft M) \cong \mathscr{B}(R, T)(M \odot N, P) \cong \mathscr{B}(R, S)(M, N \triangleright P)
\end{equation}
for all triples of 1-cells $M : R \to S$, $N : S \to T$, and $P : R \to T$.

\begin{defn}[Compare to \cref{defn:shadow}]\cite[Def 4.1]{barhite2023bicategorical}
    A \textbf{coshadow} for a closed bicategory $\B$ is a category $\mathbf{T}$ and functors \[\coshh{\text{-}}: \B(R, R) \to \mathbf{T} \] for each 0-cell $R$ of $\B$, equipped with natural isomorphisms \[\theta: \coshh{M\triangleright N} \overset{\cong}{\to} \coshh{N \triangleleft M}\] for each $M, N \in \B(R, S)$, such that the following diagrams commute whenever they make sense:

    \[\xymatrix{\coshh{(M\odot N)\triangleright P} \ar[r]^\theta_\cong \ar[d]^\cong_\theta & \coshh{P\triangleleft (M\odot N)} \ar[r]^{\coshh{t}}_\cong & \coshh{(P\triangleleft M)\triangleleft N} \ar[d]^-\theta_\cong\\
    \coshh{M\triangleright (N\triangleright P)} \ar[r]^\cong_\theta & \coshh{(N\triangleright P)\triangleleft M} \ar[r]^\cong_{\coshh{a}} & \coshh{N\triangleright (P\triangleleft M)}}\]

    \[\xymatrix{\coshh{U_R\triangleright M} \ar[r]^\theta_\cong \ar[dr]^\cong_{\coshh{\overline{r}^{-1}}} & \coshh{M\triangleleft U_R} \ar[r]^\theta_\cong \ar[d]^\cong_{\coshh{\overline{\ell}^{-1}}} & \coshh{U_R\triangleright M} \ar[dl]^{\coshh{\overline{r}^{-1}}}_\cong \\ 
    & \coshh{M} &}\]
\end{defn}

The definition of the cotrace parallels the definition of the trace in \cref{defn:bicat_trace}.

\begin{defn}\cite[Definition 4.5]{barhite2023bicategorical}
    Let $\B$ be a closed bicategory with a coshadow and $(M, M^\ast)$ a dual pair with $M\in \B(R, S)$. The \textbf{cotrace} of a 2-cell $f: M\triangleright Q \to P\triangleleft M$, denoted $\co(f)$, is the composite \small{\begin{align*}\coshh{Q} \underset{\cong}{\xto{\overline{r}}} \coshh{U_S \triangleright Q} \underset{\cong}{\xto{\coshh{\epsilon\triangleright\id}}} \coshh{(M^\ast\odot M)\triangleright Q} \underset{\cong}{\xto{\coshh{t}}} \coshh{M^\ast \triangleright (M\triangleright Q)}
    \underset{\cong}{\xto{\coshh{\id\triangleright f}}} \coshh{M^\ast\triangleright(P\triangleleft M)}\\ \hfill\underset{\cong}{\xto{\theta}} \coshh{(P\triangleleft M)\triangleleft M^\ast} \underset{\cong}{\xto{\coshh{t_\ast^{-1}}}} \coshh{P\triangleleft (M\odot M^\ast)}
    \hfill \underset{\cong}{\xto{\id\triangleleft\eta}} \coshh{P\triangleleft U_R} \underset{\cong}{\xto{\coshh{\overline{\ell}^{-1}}}} \coshh{P}\end{align*}}
\end{defn}

Not only are the definitions of the trace and cotrace parallel, but formal results for traces have almost identical proofs for cotraces.  In particular, the results in \cref{sec:remnant,sec:remnant_functors,sec:traces} have the following companions.  (We omit the proofs since they follow so closely.)

Paralleling \cref{Lambda_bicat} we have the following.

\begin{thm}\label{Lambda_bicat_coshadow}\label{coLambdaB_bicat}
    From a closed  bicategory $\mathcal{B}$ there is a bicategory $\cb$ where 
    \begin{itemize}
        \item the 0-cells are the endomorphism 1-cells of $\mathcal{B}$,
        \item the 1-cells $P\to Q$ are pairs $(M,\bar{\alpha})$ where $M$ is a dualizable 1-cell in $\mathcal{B}$ and 
        $\bar{\alpha}$ is a 2-cell 
        \[M\triangleright  P\Rightarrow Q\triangleleft M,\]
        and 
        \item the 2-cells $(M,\bar{\alpha})\Rightarrow (M',\bar{\alpha}')$ are isomorphism 2-cells $\psi\colon M\to M$ of $\mathcal{B}$ so that the diagram 
          \[\xymatrix{M\triangleright P\ar@{=>}[d]_{\psi\triangleright \id } \ar@{=>}[r]^{\bar{\alpha}} 
          & Q\triangleleft M \ar@{=>}[d]^{\id\triangleleft\psi}\\
    M'\triangleright P \ar@{=>}[r]_{\bar{\alpha}'} & Q\triangleleft M'}\]
    commutes.
    \end{itemize}
\end{thm}
We call this bicategory the \textbf{\coremant\ bicategory}.





\begin{thm}[Compare to \cref{lambda_of_a_composite}]\label{coLambdaF}\label{colambda_of_a_composite}    A strong functor $F\colon \B\to \mathcal{C}$ induces a strong functor $\cob{F}\colon \cb\to \cob{\mathcal{C}}$ of \coremant\, bicategories.
    If $F\colon \B\to \mathcal{C}$ and $G\colon \mathcal{C}\to \mathcal{D}$ are strong functors of bicategories, 
\[\cob{(F\circ G)}=(\cob{F})\circ (\cob{G}).\]
\end{thm}



\begin{thm}[Compare to \cref{thm:character}]
Let $\B$ be a closed bicategory with coshadow functors $\coshh{\text{-}}: \B(R, R) \to \mathbf{T}$. Let $\T$ be $\mathbf{T}$ regarded as a 2-category with identity 2-cells. Define $\overline{\rchi}: \cb \to \T$ by \begin{gather*} \overline{\rchi}(P) = \coshh{P},\\ \overline{\rchi}([N, \overline{\gamma}]) = \co(\overline{\gamma}),\ \text{and} \\ \overline{\rchi}(\lone{\psi}) = \id_{\co(\overline{\gamma})}, \end{gather*} where $\co$ is the bicategorical cotrace with coshadow $\coshh{\text{-}}$. Then $\overline{\rchi}$ is a strict 2-functor.
\end{thm}
 

 \begin{defn}[Compare to \cref{lab:def:character}]\label{cocharacter}
    Let $\B$ be a closed bicategory with coshadow $\coshh{}$ valued in $\mathbf{T}$ and $\rho\colon BG\to \B$ be a 2-representation. Let $\T$ be $\mathbf{T}$ regarded as a 2-category with identity 2-cells. The \textbf{cocharacter} of $\rho$ is the composite
    \[\Lambda BG\xto{\cob{\rho}} \cob{\B} \xto{\overline{\chi}}\T,\] where $\cob{\rho}(g) = \rho(g)$ and $\cob{\rho} = \overline{\gamma}$.
\end{defn}

\begin{ex}\label{ex:GK}
    Let $\rho: BG \to \B$ be a 2-representation, where $\B$ is a closed bicategory. Let the coshadow in the definition of $\overline{\rchi}$ be $\Tr$, the categorical trace as defined in \cite[3.1]{GK}, with target category $T$. The categorical character, $\Tr(\rho)$, from \cite[4.8]{GK} is given by the commutative diagram: \[\xymatrix{\Lambda BG \ar[rr]^{\Tr(\rho)} \ar[dr]_{\overline{\Lambda}\rho} & & \T\\
    & \cb \ar[ur]_{\overline{\rchi}} &}\] On the level of 1-cells, $\overline{\rchi}(\overline{\gamma}) = \co(\overline{\gamma})$. 
\end{ex}


Just as \cref{lambda_of_a_composite} implies that character commutes with restriction, \cref{colambda_of_a_composite} implies that the cocharacter commutes with restriction. 

